% ============================================================================
% ARTIGO CIENTÍFICO
% Sistema Autônomo de Autenticação Contextual com Blockchain e
% Inteligência Artificial
% ============================================================================
% Autor: João Guedes de Brito
% Orientador: Cleber Jorge Lira de Santana
% Instituição: IFBA - Programa de Pós-Graduação em Engenharia de Sistemas e Produtos
% ============================================================================

\documentclass[12pt,a4paper]{article}
\usepackage[utf8]{inputenc}
\usepackage[T1]{fontenc}
\usepackage[brazil]{babel}
\usepackage{graphicx}
\usepackage{amsmath}
\usepackage{hyperref}
\usepackage{booktabs}
\usepackage{float}
\usepackage{geometry}
\usepackage{indentfirst}
\usepackage{setspace}
\usepackage{cite}
\usepackage{url}
\usepackage{listings}
\usepackage{xcolor}

\geometry{left=3cm, right=2cm, top=3cm, bottom=2cm}
\onehalfspacing

\title{\textbf{Sistema Autônomo de Autenticação Contextual com Blockchain e Inteligência Artificial: Uma Abordagem Baseada em Microserviços para Segurança Adaptativa em Aplicações Web}}

\author{
  João Guedes de Brito\textsuperscript{1}, Cleber Jorge Lira de Santana\textsuperscript{1} \\
  \small \textsuperscript{1}Programa de Pós-Graduação em Engenharia de Sistemas e Produtos \\
  \small Instituto Federal da Bahia (IFBA) -- Salvador, BA -- Brasil \\
  \small \texttt{joao.brito@ifba.edu.br, cleber.santana@ifba.edu.br}
}

\date{}

\begin{document}

\maketitle

% ============================================================================
% RESUMO
% ============================================================================
\begin{abstract}
Este artigo apresenta o desenvolvimento e a avaliação de um sistema autônomo de autenticação contextual que integra Inteligência Artificial (IA) e tecnologia Blockchain para reforçar a segurança de acesso em aplicações web. A arquitetura proposta baseia-se em microserviços e articula cinco camadas funcionais: aplicação modular, integração contínua, segurança automatizada, inteligência artificial e registro imutável. O módulo de IA emprega um modelo ensemble (Random Forest, Decision Tree e Logistic Regression) para classificação de risco em tempo real, analisando dez atributos contextuais -- incluindo geolocalização, dispositivo, horário e padrões comportamentais do usuário. A camada Blockchain garante imutabilidade e rastreabilidade forense dos logs de autenticação, atendendo a requisitos normativos como a LGPD. Os resultados experimentais demonstram acurácia de 91\% na detecção de acessos anômalos, com latência média de 320ms no fluxo completo de autenticação contextual e integridade de 100\% nas transações registradas em Blockchain. O sistema foi implementado com Spring Boot, Angular, Keycloak, Scikit-learn e Hyperledger/Ethereum, orquestrado via Docker, com monitoramento em tempo real através do ELK Stack e Grafana. A solução proposta demonstra viabilidade técnica como alternativa adaptativa e auditável aos métodos tradicionais de autenticação.

\textbf{Palavras-chave:} Autenticação Contextual; Inteligência Artificial; Blockchain; Segurança da Informação; Machine Learning; Microserviços; LGPD.
\end{abstract}

\vspace{0.5cm}

\begin{center}
\textbf{ABSTRACT}
\end{center}

\noindent This paper presents the development and evaluation of an autonomous contextual authentication system that integrates Artificial Intelligence (AI) and Blockchain technology to enhance access security in web applications. The proposed architecture is based on microservices and comprises five functional layers: modular application, continuous integration, automated security, artificial intelligence, and immutable logging. The AI module employs an ensemble model (Random Forest, Decision Tree, and Logistic Regression) for real-time risk classification, analyzing ten contextual attributes -- including geolocation, device, time, and user behavioral patterns. The Blockchain layer ensures immutability and forensic traceability of authentication logs, meeting regulatory requirements such as Brazil's LGPD. Experimental results demonstrate 91\% accuracy in anomalous access detection, with an average latency of 320ms for the complete contextual authentication flow and 100\% integrity in Blockchain-recorded transactions.

\noindent \textbf{Keywords:} Contextual Authentication; Artificial Intelligence; Blockchain; Information Security; Machine Learning; Microservices; LGPD.

% ============================================================================
% 1. INTRODUÇÃO
% ============================================================================
\section{Introdução}

Em um cenário digital marcado pela crescente sofisticação dos ataques cibernéticos e pela expansão dos serviços baseados em web, a autenticação de usuários permanece como um dos pilares mais críticos da segurança da informação. Métodos tradicionais, como senhas fixas e autenticação em dois fatores (2FA), embora amplamente adotados, apresentam limitações significativas diante de ataques como \textit{credential stuffing}, \textit{session hijacking} e engenharia social \cite{owasp2023}. Essas abordagens, predominantemente estáticas, não consideram o contexto comportamental e ambiental do usuário no momento do acesso, criando lacunas exploráveis por agentes maliciosos.

A crescente demanda por conformidade regulatória, como a Lei Geral de Proteção de Dados (LGPD) no Brasil \cite{brasil2018}, e as diretrizes internacionais do National Institute of Standards and Technology (NIST) \cite{nist2020}, reforçam a necessidade de sistemas de autenticação que combinem adaptabilidade, auditabilidade e transparência. Nesse contexto, a convergência entre Inteligência Artificial (IA), tecnologia Blockchain e arquiteturas distribuídas emerge como uma abordagem promissora para superar as limitações dos mecanismos convencionais.

Este trabalho propõe e implementa um sistema autônomo de autenticação contextual que integra: (i) análise de risco em tempo real baseada em modelos supervisionados de Machine Learning; (ii) registro imutável de eventos de autenticação via Blockchain; e (iii) uma arquitetura de microserviços orientada por práticas DevSecOps. O sistema analisa dez atributos contextuais por acesso -- incluindo localização geográfica, dispositivo, horário, padrão comportamental e indicadores de rede -- para tomar decisões autônomas sobre a concessão, monitoramento, exigência de autenticação multifator ou bloqueio de acesso.

A principal contribuição deste trabalho reside na integração sistêmica de três camadas de segurança complementares -- autenticação baseada em identidade (Keycloak/OAuth2), análise preditiva de risco (IA) e auditoria imutável (Blockchain) -- em uma arquitetura modular, escalável e reprodutível. Os resultados experimentais demonstram a viabilidade técnica da abordagem, com acurácia de 91\% na classificação de acessos anômalos e latência compatível com ambientes de produção.

O restante deste artigo está organizado da seguinte forma: a Seção 2 apresenta a fundamentação teórica; a Seção 3 descreve a arquitetura e a implementação do sistema; a Seção 4 detalha a metodologia experimental; a Seção 5 apresenta e discute os resultados; e a Seção 6 conclui o trabalho com considerações finais e direções futuras.

% ============================================================================
% 2. FUNDAMENTAÇÃO TEÓRICA
% ============================================================================
\section{Fundamentação Teórica}

\subsection{Autenticação Contextual e Segurança Adaptativa}

A autenticação contextual representa uma evolução dos modelos tradicionais ao incorporar variáveis ambientais e comportamentais no processo decisório de concessão de acesso. Diferentemente da autenticação baseada exclusivamente em credenciais (\textit{something you know}), a abordagem contextual analisa fatores como localização geográfica, dispositivo utilizado, horário de acesso e histórico comportamental do usuário para calcular um nível de risco associado à tentativa de autenticação \cite{nist2020}.

O \textit{Secure Software Development Framework} (SSDF), publicado pelo NIST \cite{nist2022}, estabelece diretrizes para o desenvolvimento seguro que incluem automação de controles, rastreabilidade de alterações e integração de mecanismos de verificação contínua. A abordagem \textit{Secure by Design}, incentivada pela Cybersecurity and Infrastructure Security Agency (CISA) \cite{cisa2023}, complementa essa visão ao propor que a segurança seja considerada desde a concepção do sistema, não como camada adicional posterior.

\subsection{Inteligência Artificial Aplicada à Cibersegurança}

A aplicação de técnicas de Machine Learning na cibersegurança tem demonstrado resultados significativos na detecção de anomalias e classificação de ameaças. Modelos supervisionados permitem classificar comportamentos conhecidos com base em dados históricos rotulados, enquanto modelos não supervisionados possibilitam a identificação de padrões anômalos previamente não catalogados \cite{goodfellow2016}.

O \textit{AI Risk Management Framework} (AI RMF 1.0) do NIST \cite{nist2023ai} estabelece que sistemas baseados em IA devem garantir confiabilidade, rastreabilidade e governança adequada, implicando registro auditável de decisões automatizadas e monitoramento contínuo do desempenho dos modelos. Métricas como precisão, \textit{recall} e F1-\textit{score} são fundamentais para a avaliação rigorosa dos classificadores, garantindo equilíbrio entre sensibilidade e especificidade.

Modelos \textit{ensemble}, que combinam múltiplos algoritmos de classificação por meio de mecanismos de votação, têm se mostrado particularmente eficazes em tarefas de detecção de anomalias, pois mitigam vieses individuais de cada classificador e aumentam a robustez preditiva do sistema \cite{scikit2024}.

\subsection{Blockchain como Mecanismo de Integridade e Auditoria}

A tecnologia Blockchain, fundamentada no encadeamento criptográfico de blocos descrito por Nakamoto \cite{nakamoto2008}, oferece garantias de imutabilidade, rastreabilidade e não repúdio para registros digitais. No contexto de auditoria de segurança, blockchains permissionadas são adequadas para ambientes corporativos e governamentais, pois combinam controle de acesso institucional com as propriedades de integridade distribuída \cite{nist2018blockchain}.

Conforme recomendado pelo NIST \cite{nist2020}, a proteção da integridade de logs e trilhas de auditoria é requisito essencial em sistemas de segurança. A utilização de Blockchain para registrar \textit{hashes} criptográficos de eventos de autenticação garante que nenhum registro possa ser alterado retroativamente, oferecendo evidências forenses confiáveis e conformidade com regulamentações como a LGPD.

\subsection{Arquitetura de Microserviços e DevSecOps}

A adoção do paradigma de microserviços, conforme descrito por Newman \cite{newman2021}, favorece a decomposição funcional, o isolamento de domínios e a escalabilidade independente de componentes. Sob a perspectiva da segurança, essa abordagem permite a aplicação de políticas granulares de autenticação, autorização e monitoramento para cada serviço.

O DevSecOps integra controles de segurança desde as fases iniciais do ciclo de desenvolvimento, promovendo a abordagem \textit{Shift Left Security}. A incorporação de análise estática (SAST), análise dinâmica (DAST) e escaneamento de dependências reduz significativamente a exposição a vulnerabilidades críticas, especialmente aquelas listadas no OWASP Top 10 \cite{owasp2023}.

% ============================================================================
% 3. ARQUITETURA E IMPLEMENTAÇÃO
% ============================================================================
\section{Arquitetura e Implementação}

\subsection{Visão Geral da Arquitetura}

A arquitetura proposta organiza-se em cinco camadas integradas, conforme ilustrado na Figura \ref{fig:arquitetura}:

\begin{enumerate}
    \item \textbf{Camada de Aplicação Modular} -- composta por serviços independentes (backend Spring Boot, frontend Angular) e APIs autenticadas via JWT/OAuth2, com Keycloak como provedor de identidade.
    \item \textbf{Camada de Integração Contínua} -- responsável pela orquestração automatizada do pipeline DevSecOps, incluindo testes unitários, de integração e de segurança.
    \item \textbf{Camada de Segurança Automatizada} -- executa verificações SAST, DAST e análise de dependências contra bases de vulnerabilidades conhecidas.
    \item \textbf{Camada de Inteligência Artificial} -- classifica riscos e detecta anomalias com base em modelos supervisionados de aprendizado de máquina.
    \item \textbf{Camada de Registro Imutável} -- assegura rastreabilidade e integridade forense por meio de Blockchain permissionada.
\end{enumerate}

% \begin{figure}[H]
%     \centering
%     \includegraphics[width=0.9\textwidth]{figuras/arquitetura-cinco-camadas.png}
%     \caption{Arquitetura em Cinco Camadas do Sistema Proposto}
%     \label{fig:arquitetura}
% \end{figure}

\subsection{Módulo de Autenticação Contextual}

O módulo de autenticação contextual opera como intermediário inteligente no fluxo de login. Quando um usuário inicia uma tentativa de autenticação, o sistema captura automaticamente dez atributos contextuais, detalhados na Tabela \ref{tab:features}.

\begin{table}[H]
\centering
\caption{Atributos contextuais extraídos por tentativa de autenticação}
\label{tab:features}
\begin{tabular}{@{}lll@{}}
\toprule
\textbf{\#} & \textbf{Atributo} & \textbf{Descrição} \\
\midrule
1  & \texttt{hour\_of\_day}              & Hora do acesso (0--23) \\
2  & \texttt{day\_of\_week}              & Dia da semana (0--6) \\
3  & \texttt{login\_attempts\_last\_hour} & Tentativas na última hora \\
4  & \texttt{is\_new\_device}            & Dispositivo não reconhecido (0/1) \\
5  & \texttt{is\_new\_location}          & Localização inédita (0/1) \\
6  & \texttt{is\_vpn}                    & Uso de VPN detectado (0/1) \\
7  & \texttt{is\_tor}                    & Uso de rede Tor detectado (0/1) \\
8  & \texttt{session\_duration\_avg}     & Duração média de sessão (s) \\
9  & \texttt{pages\_per\_session\_avg}   & Páginas por sessão (média) \\
10 & \texttt{time\_since\_last\_login}   & Tempo desde último login (h) \\
\bottomrule
\end{tabular}
\end{table}

O backend (Spring Boot) implementa cinco estratégias de análise de risco, seguindo o padrão \textit{Strategy}: análise comportamental, de dispositivo, de localização, de rede e temporal. Cada estratégia contribui com uma pontuação parcial que alimenta o modelo de IA para a classificação final.

\subsection{Modelo de Machine Learning}

O serviço de IA foi implementado como microsserviço independente em Python com FastAPI, expondo endpoints RESTful para predição (\texttt{/predict}), treinamento (\texttt{/train}) e verificação de saúde (\texttt{/health}).

O modelo de classificação utiliza um \textit{ensemble} baseado em \textit{Soft Voting}, combinando três algoritmos:

\begin{itemize}
    \item \textbf{Random Forest} -- 100 estimadores, profundidade máxima 10;
    \item \textbf{Decision Tree} -- profundidade máxima 8;
    \item \textbf{Logistic Regression} -- 500 iterações, otimizador \textit{lbfgs}.
\end{itemize}

O mecanismo de \textit{Soft Voting} calcula a média ponderada das probabilidades preditas por cada classificador, resultando em um \textit{score} de risco contínuo no intervalo $[0, 1]$. Os dados de entrada são normalizados via \texttt{StandardScaler} antes da inferência. O modelo foi treinado com 5.000 amostras sintéticas geradas a partir de regras heurísticas que simulam cenários realistas de autenticação.

A classificação de risco segue quatro níveis, conforme a Tabela \ref{tab:risk_levels}.

\begin{table}[H]
\centering
\caption{Níveis de risco e decisões automatizadas}
\label{tab:risk_levels}
\begin{tabular}{@{}llll@{}}
\toprule
\textbf{Nível} & \textbf{Score} & \textbf{Decisão} & \textbf{Ação} \\
\midrule
Baixo    & $< 0{,}3$          & PERMITIR                   & Acesso liberado \\
Médio    & $[0{,}3; 0{,}6)$   & PERMITIR\_COM\_MONITORAMENTO & Acesso com alerta \\
Alto     & $[0{,}6; 0{,}8)$   & EXIGIR\_MFA                & Autenticação multifator \\
Crítico  & $\geq 0{,}8$       & BLOQUEAR                   & Acesso negado \\
\bottomrule
\end{tabular}
\end{table}

\subsection{Camada Blockchain}

A camada de registro imutável utiliza duas implementações complementares:

\begin{itemize}
    \item \textbf{Ethereum} (via Web3j 4.10.3) -- para registro de \textit{hashes} criptográficos de eventos de autenticação em rede pública ou privada;
    \item \textbf{Hyperledger Fabric} -- para ambientes corporativos que requerem governança permissionada e controle institucional.
\end{itemize}

Cada evento de autenticação gera uma transação contendo: tipo do evento (sucesso, negação, MFA requerido), identificador do usuário (anonimizado via \textit{hash} para conformidade com LGPD), endereço IP, \textit{fingerprint} do dispositivo, \textit{score} de risco, decisão tomada e \textit{timestamp}. A integridade dos registros é verificável por meio de endpoint específico que realiza validação criptográfica da cadeia.

\subsection{Gerenciamento de Identidade}

O Keycloak foi configurado como provedor central de identidade, suportando os protocolos OAuth 2.0 e OpenID Connect. A configuração inclui:

\begin{itemize}
    \item \textbf{Proteção contra força bruta}: limite de 5 tentativas falhas, bloqueio temporário de 900 segundos;
    \item \textbf{Clientes OIDC}: cliente público para o frontend Angular (com PKCE) e cliente confidencial para o backend (\textit{bearer-only});
    \item \textbf{Tokens}: \textit{access token} com validade de 5 minutos, SSO com \textit{idle timeout} de 30 minutos.
\end{itemize}

\subsection{Observabilidade e Monitoramento}

O sistema integra o ELK Stack (Elasticsearch, Logstash e Kibana) e Grafana para monitoramento em tempo real. O pipeline Logstash classifica eventos por tipo (autenticação, análise de risco, blockchain, alertas) e indexa no Elasticsearch com padrão temporal (\texttt{auth-logs-YYYY.MM.dd}). Dashboards Grafana permitem visualização de métricas de segurança, incluindo taxa de bloqueios, distribuição de \textit{scores} de risco e latência dos serviços.

\subsection{Stack Tecnológica}

A Tabela \ref{tab:tech_stack} sumariza as principais tecnologias empregadas.

\begin{table}[H]
\centering
\caption{Stack tecnológica do sistema}
\label{tab:tech_stack}
\begin{tabular}{@{}lll@{}}
\toprule
\textbf{Camada} & \textbf{Tecnologia} & \textbf{Versão} \\
\midrule
Backend           & Spring Boot        & 3.2.3 \\
Frontend          & Angular            & 17.2.0 \\
Identidade        & Keycloak           & 24.0 \\
ML/IA             & Scikit-learn       & 1.4.0 \\
API de IA         & FastAPI            & 0.109.0 \\
Blockchain        & Web3j (Ethereum)   & 4.10.3 \\
Blockchain        & Hyperledger Fabric & 2.5 \\
Banco de Dados    & PostgreSQL         & 15 \\
Cache             & Redis              & 7 \\
Containerização   & Docker Compose     & Latest \\
Logs              & ELK Stack          & 8.12.0 \\
Monitoramento     & Grafana            & 10.3.1 \\
\bottomrule
\end{tabular}
\end{table}

% ============================================================================
% 4. METODOLOGIA EXPERIMENTAL
% ============================================================================
\section{Metodologia Experimental}

\subsection{Configuração do Ambiente}

O ambiente experimental foi construído com infraestrutura containerizada via Docker, com cada componente implantado como microsserviço independente em uma rede bridge compartilhada (\texttt{auth-network}). Os dez serviços orquestrados incluem: PostgreSQL, Redis, Keycloak, backend Spring Boot, frontend Angular, serviço de IA (FastAPI), Elasticsearch, Logstash, Kibana e Grafana.

\subsection{Eixos de Validação}

Os experimentos foram estruturados em três eixos:

\begin{enumerate}
    \item \textbf{Avaliação do modelo de ML}: classificação de acessos legítimos e anômalos com validação cruzada ($k$-\textit{fold}, $k=10$);
    \item \textbf{Cadeia de autenticação contextual}: integração Keycloak $\rightarrow$ captura de contexto $\rightarrow$ classificação por IA $\rightarrow$ decisão $\rightarrow$ registro em Blockchain;
    \item \textbf{Integridade Blockchain}: verificação de imutabilidade e consistência criptográfica das transações.
\end{enumerate}

\subsection{Conjunto de Dados}

O modelo foi treinado com 5.000 amostras sintéticas de logs de autenticação, contendo os dez atributos contextuais descritos na Tabela \ref{tab:features}. Os dados foram rotulados em quatro categorias de risco (baixo, médio, alto, crítico) com base em regras heurísticas derivadas de cenários reais de ataque. A distribuição foi balanceada para mitigar viés de classe.

\subsection{Cenários de Teste}

Foram simulados cenários controlados representativos de ameaças reais:

\begin{itemize}
    \item Acesso legítimo a partir de dispositivo e localização habituais;
    \item Acesso a partir de novo dispositivo em localização conhecida;
    \item Acesso via VPN/Tor em horário atípico (2h--6h);
    \item Múltiplas tentativas falhas seguidas de tentativa com credenciais válidas (\textit{credential stuffing});
    \item Acesso após período prolongado de inatividade ($>30$ dias).
\end{itemize}

% ============================================================================
% 5. RESULTADOS E DISCUSSÃO
% ============================================================================
\section{Resultados e Discussão}

\subsection{Desempenho do Modelo de Machine Learning}

Os resultados da avaliação do modelo \textit{ensemble} são apresentados na Tabela \ref{tab:ml_results}.

\begin{table}[H]
\centering
\caption{Métricas de desempenho do modelo de classificação}
\label{tab:ml_results}
\begin{tabular}{@{}ll@{}}
\toprule
\textbf{Métrica} & \textbf{Valor} \\
\midrule
Acurácia            & 91,0\% \\
Precisão            & 89,4\% \\
Recall              & 93,1\% \\
F1-Score            & 91,2\% \\
Taxa de Falsos Positivos & 8,6\% \\
\bottomrule
\end{tabular}
\end{table}

O F1-\textit{score} de 91,2\% indica equilíbrio adequado entre precisão e sensibilidade. O \textit{recall} de 93,1\% é particularmente relevante no contexto de segurança, pois prioriza a detecção de acessos genuinamente anômalos, minimizando falsos negativos que representariam acessos maliciosos não detectados. A taxa de falsos positivos de 8,6\%, embora aceitável para a fase atual de validação com dados sintéticos, representa uma oportunidade de otimização com a incorporação futura de dados reais de produção.

\subsection{Latência da Autenticação Contextual}

A latência média do fluxo completo de autenticação contextual foi medida em 320ms, conforme decomposição na Tabela \ref{tab:latency}.

\begin{table}[H]
\centering
\caption{Decomposição da latência do fluxo de autenticação}
\label{tab:latency}
\begin{tabular}{@{}ll@{}}
\toprule
\textbf{Etapa} & \textbf{Latência Média} \\
\midrule
Validação de credenciais (Keycloak)      & 80--100ms \\
Captura de contexto                       & 20--30ms \\
Classificação de risco (IA)               & 120--180ms \\
Registro em Blockchain                    & 40--60ms \\
\midrule
\textbf{Total}                            & \textbf{$\approx$320ms} \\
\bottomrule
\end{tabular}
\end{table}

A latência adicional introduzida pelo módulo de IA (120--180ms) é o principal gargalo identificado. Estratégias de mitigação incluem cache de decisões para usuários com padrões reconhecidos e otimização do modelo com técnicas de quantização. Não obstante, o tempo total de 320ms é compatível com padrões de experiência do usuário em aplicações web, que tipicamente toleram latências de até 500ms em processos de autenticação.

\subsection{Integridade Blockchain}

Os testes de integridade demonstraram consistência de 100\% nas transações registradas. A validação criptográfica da cadeia de blocos confirmou que nenhum registro foi alterado ou corrompido durante os ciclos de teste. Cada evento de autenticação gerou uma transação verificável, com \textit{hash} de confirmação e metadados completos.

A utilização de anonimização por \textit{hash} dos dados pessoais antes do registro em Blockchain atende ao princípio de minimização de dados da LGPD, garantindo rastreabilidade forense sem exposição de informações pessoais identificáveis.

\subsection{Análise Comparativa}

A Tabela \ref{tab:comparison} apresenta uma comparação qualitativa entre o sistema proposto e abordagens tradicionais de autenticação.

\begin{table}[H]
\centering
\caption{Comparação entre abordagens de autenticação}
\label{tab:comparison}
\begin{tabular}{@{}lccc@{}}
\toprule
\textbf{Característica} & \textbf{Senha} & \textbf{2FA} & \textbf{Proposta} \\
\midrule
Análise contextual       & Não & Não & Sim \\
Adaptabilidade           & Não & Não & Sim \\
Auditoria imutável       & Não & Não & Sim \\
Detecção de anomalias    & Não & Parcial & Sim \\
Conformidade LGPD        & Parcial & Parcial & Sim \\
Resistência a \textit{credential stuffing} & Não & Parcial & Sim \\
Decisão autônoma         & Não & Não & Sim \\
\bottomrule
\end{tabular}
\end{table}

\subsection{Discussão}

Os resultados confirmam a hipótese de que a integração de IA contextual e Blockchain ao processo de autenticação eleva significativamente o nível de segurança em comparação com métodos estáticos. A abordagem adaptativa identificou padrões suspeitos em cenários de \textit{credential stuffing} e \textit{session hijacking} que seriam invisíveis a sistemas convencionais.

A arquitetura de microserviços demonstrou adequação para o cenário proposto, permitindo escalabilidade independente dos componentes e facilitando a evolução incremental do sistema. O Keycloak, como camada central de identidade, ofereceu suporte nativo a OAuth2 e OpenID Connect sem desenvolvimento adicional de infraestrutura.

As principais limitações identificadas incluem: (i) a utilização de dados sintéticos para treinamento, que impacta a generalização do modelo; (ii) a complexidade operacional da rede Hyperledger Fabric em ambientes de desenvolvimento; e (iii) a necessidade de definição cuidadosa de políticas de minimização e anonimização de dados pessoais para plena conformidade com a LGPD.

% ============================================================================
% 6. CONCLUSÃO
% ============================================================================
\section{Conclusão}

Este trabalho apresentou o desenvolvimento e a avaliação experimental de um sistema autônomo de autenticação contextual que integra Inteligência Artificial e Blockchain em uma arquitetura de microserviços orientada por práticas DevSecOps. Os resultados demonstram a viabilidade técnica da abordagem, com acurácia de 91\% na classificação de risco, latência compatível com ambientes de produção e integridade total dos registros em Blockchain.

A contribuição principal reside na integração sistêmica de três camadas complementares de segurança -- identidade, inteligência preditiva e auditoria imutável -- em um sistema modular e reprodutível. A arquitetura proposta oferece um novo paradigma para autenticação em ambientes críticos, superando as limitações reativas dos métodos tradicionais por meio de uma abordagem proativa, autônoma e orientada por evidências.

Como trabalhos futuros, destacam-se: (i) a coleta e utilização de dados reais de autenticação em ambientes controlados para retreinamento do modelo; (ii) a incorporação de técnicas de aprendizado contínuo (\textit{online learning}) para adaptação dinâmica aos padrões dos usuários; (iii) a avaliação de desempenho em cenários de alta concorrência com testes de carga; (iv) a extensão da solução para contextos de IoT e autenticação federada; e (v) a implementação de mecanismos de explicabilidade (XAI) para as decisões do modelo de IA.

% ============================================================================
% REFERÊNCIAS
% ============================================================================
\bibliographystyle{abntex2-num}

\begin{thebibliography}{99}

\bibitem{brasil2018}
BRASIL. \textbf{Lei nº 13.709, de 14 de agosto de 2018}. Lei Geral de Proteção de Dados Pessoais (LGPD). Diário Oficial da União, Brasília, DF, 15 ago. 2018.

\bibitem{cisa2023}
CISA -- CYBERSECURITY AND INFRASTRUCTURE SECURITY AGENCY. \textbf{Secure by Design Initiative}. Washington, DC: CISA, 2023.

\bibitem{eua2025}
ESTADOS UNIDOS. \textbf{Executive Order 14144}: Strengthening and Promoting Innovation in the Nation's Cybersecurity. Washington, DC: The White House, 2025.

\bibitem{goodfellow2016}
GOODFELLOW, I.; BENGIO, Y.; COURVILLE, A. \textbf{Deep Learning}. Cambridge: MIT Press, 2016.

\bibitem{hyperledger2024}
HYPERLEDGER FOUNDATION. \textbf{Hyperledger Fabric Documentation}. Versão 2.5. Linux Foundation, 2024.

\bibitem{keycloak2024}
KEYCLOAK. \textbf{Keycloak Documentation: Server Administration Guide}. Versão 24. Red Hat, 2024.

\bibitem{nakamoto2008}
NAKAMOTO, S. \textbf{Bitcoin}: A Peer-to-Peer Electronic Cash System. 2008.

\bibitem{nist2023ai}
NATIONAL INSTITUTE OF STANDARDS AND TECHNOLOGY (NIST). \textbf{AI Risk Management Framework} (AI RMF 1.0). NIST AI 100-1. Gaithersburg: NIST, 2023.

\bibitem{nist2018blockchain}
NATIONAL INSTITUTE OF STANDARDS AND TECHNOLOGY (NIST). \textbf{Blockchain Technology Overview}. NISTIR 8202. Gaithersburg: NIST, 2018.

\bibitem{nist2022}
NATIONAL INSTITUTE OF STANDARDS AND TECHNOLOGY (NIST). \textbf{Secure Software Development Framework} (SSDF) Version 1.1: SP 800-218. Gaithersburg: NIST, 2022.

\bibitem{nist2020}
NATIONAL INSTITUTE OF STANDARDS AND TECHNOLOGY (NIST). \textbf{Security and Privacy Controls for Information Systems and Organizations}. SP 800-53 Rev. 5. Gaithersburg: NIST, 2020.

\bibitem{newman2021}
NEWMAN, S. \textbf{Building Microservices}: Designing Fine-Grained Systems. 2. ed. Sebastopol: O'Reilly Media, 2021.

\bibitem{owasp2023}
OWASP FOUNDATION. \textbf{OWASP Top 10} -- 2023 Edition. 2023.

\bibitem{scikit2024}
SCIKIT-LEARN DEVELOPERS. \textbf{Scikit-learn}: Machine Learning in Python. Versão 1.4. 2024.

\end{thebibliography}

\end{document}
